\section*{अध्याय \ref{ch:g12:buffers-predominance} --- बफ़र और प्रभाविता आरेख}

\begin{solution}{exo:g12:buffers-predominance:1}
\omterm{def:g9:acids-bases-ph:ph}{pH} 2: \ce{CH3COOH}। \omterm{def:g9:acids-bases-ph:ph}{pH} 4.76: दोनों समान मात्राओं में। \omterm{def:g9:acids-bases-ph:ph}{pH} 7: \ce{CH3COO-}।
\end{solution}

\begin{solution}{exo:g12:buffers-predominance:2}
\omterm{def:g9:acids-bases-ph:ph}{pH} 4: $1/(1 + 10^{0.76}) = 0.15$। \omterm{def:g9:acids-bases-ph:ph}{pH} 6: $1/(1 + 10^{-1.24}) = 0.95$।
\end{solution}

\begin{solution}{exo:g12:buffers-predominance:3}
पीला; हरा (अपने परास के भीतर); नीला।
\end{solution}

\begin{solution}{exo:g12:buffers-predominance:4}
\omterm{def:g9:acids-bases-ph:ph}{pH} 7.4 पर, जो 9.26 से नीचे है: \ce{NH4+}। \omterm{def:g9:acids-bases-ph:ph}{pH} 11 पर: \ce{NH3}।
\end{solution}

\begin{solution}{exo:g12:buffers-predominance:5}
सूचक स्वयं एक \omterm{def:g12:ka-and-pka:bronsted}{अम्ल--क्षारक युग्म} है: बड़ी मात्रा में वह उसी \omterm{def:g9:acids-bases-ph:ph}{pH} को
बदल देता जिसे उसे दिखाना है।
\end{solution}

\begin{solution}{exo:g12:buffers-predominance:6}
$10^{\mathrm{pH} - 4.76}$: 0.1; 1; 10।
\end{solution}

\begin{solution}{exo:g12:buffers-predominance:7}
फ़ीनॉल्फ़्थेलीन (8.0--10.0) या थाइमॉल ब्लू (8.0--9.1): उनके परासों में
8.7 आता है। मेथिल रेड 4.2 और 6.3 के बीच बदलता है, 8.7 से बहुत दूर: वह
बहुत पहले ही बदल चुका होता।
\end{solution}

\begin{solution}{exo:g12:buffers-predominance:8}
$\mathrm{pH} = 4.76 + \log(0.10/0.20) = 4.76 - 0.30 = 4.46$।
\end{solution}

\begin{solution}{exo:g12:buffers-predominance:9}
\omterm{def:g9:acids-bases-ph:ph}{pH} 1 पर प्रभावी रूप \omterm{def:g9:ions:ion}{धनायन} \ce{H3N+-CH2-COOH} है, जिसका आवेश $+1$ है। \omterm{def:g9:acids-bases-ph:ph}{pH} 6 पर ज़्विटर \omterm{def:g9:ions:ion}{आयन},
जिसका कुल आवेश 0 है। \omterm{def:g9:acids-bases-ph:ph}{pH} 12 पर \omterm{def:g9:ions:ion}{ऋणायन} \ce{H2N-CH2-COO-} है, जिसका आवेश $-1$ है।
\end{solution}

\begin{solution}{exo:g12:buffers-predominance:10}
पानी: 7.0 से $-\log(\num{1.0e-4}/0.101) = 3.0$ तक, 4 इकाइयों की गिरावट।
बफ़र: 4.76 से $4.76 + \log(9.9/10.1) = 4.75$ तक, 0.01 की गिरावट।
\end{solution}

\begin{solution}{exo:g12:buffers-predominance:11}
\ce{NH4+}/\ce{NH3} ($pK_a = 9.26$)। अनुपात
$[\ce{NH3}]/[\ce{NH4+}] = 10^{9.5 - 9.26} = 1.7$।
\end{solution}

\begin{solution}{exo:g12:buffers-predominance:12}
\omterm{def:g9:acids-bases-ph:ph}{pH} एक इकाई तब गिरता है जब अनुपात $1/10$ तक पहुँचता है:
$(10 - n)/(10 + n) = 0.1$, इसलिए $n = \qty{8.2}{mmol}$। जिस बफ़र में हर रूप
अधिक हो, वह अनुपात में उसी परिवर्तन के लिए अधिक अम्ल या क्षारक सोख सकता है:
उसकी धारिता अधिक होती है।
\end{solution}

\begin{solution}{exo:g12:buffers-predominance:13}
अनुपात $10^{5.0 - 4.76} = 1.74$; अम्ल \qty{10}{mmol} है, इसलिए
\qty{17.4}{mmol} एथेनोएट, अर्थात विलयन के \qty{174}{mL}।
\end{solution}

\begin{solution}{exo:g12:buffers-predominance:14}
इसके दो \omterm{def:g12:ka-and-pka:bronsted}{अम्ल--क्षारक युग्म} हैं, दो $pK_a$ के साथ, इसलिए रंग के दो परिवर्तन:
\omterm{def:g9:acids-bases-ph:ph}{pH} 1 पर लाल, \omterm{def:g9:acids-bases-ph:ph}{pH} 5 पर पीला, \omterm{def:g9:acids-bases-ph:ph}{pH} 10 पर नीला।
\end{solution}

\begin{solution}{exo:g12:buffers-predominance:15}
\omterm{def:g10:concentration-and-dilution:dilution}{तनुकरण} दोनों सांद्रताओं को 10 से विभाजित करता है: अनुपात, और इसलिए \omterm{def:g9:acids-bases-ph:ph}{pH}
(4.76), नहीं बदलते। दस गुना तनु किए गए \omterm{def:g12:ka-and-pka:strong-acid}{प्रबल अम्ल} का \omterm{def:g9:acids-bases-ph:ph}{pH} एक
इकाई बढ़ जाता है।
\end{solution}

\begin{solution}{pb:g12:buffers-predominance:1}
\textbf{1.} \ce{CO2 + 2H2O <=> HCO3- + H3O+}।

\textbf{2.} घुली हुई कार्बन डाइऑक्साइड (पानी के साथ) अम्ल है,
हाइड्रोजनकार्बोनेट क्षारक।

\textbf{3.} सारणी का मान शुद्ध पानी में \qty{25}{\celsius} के लिए है;
रक्त \qty{37}{\celsius} पर है और अन्य \omterm{def:g9:ions:ion}{आयनों} से भरा है, जो
स्थिरांक को बदलते हैं (और शरीर-क्रिया विज्ञानी सारी घुली हुई कार्बन डाइऑक्साइड गिनते हैं)।

\textbf{4.} \qty{24}{mmol/L} पर, $0.024 \times 5.0 = \qty{0.12}{mol}$।

\textbf{5.} 6.1 पर कटा एक \omterm{def:g9:acids-bases-ph:ph}{pH} अक्ष: नीचे \ce{CO2}, ऊपर \ce{HCO3-}।

\textbf{6.} 7.4 पर, जो 6.1 से ऊपर है: हाइड्रोजनकार्बोनेट।

\textbf{7.} थोड़ा क्षारकीय: उसका \omterm{def:g9:acids-bases-ph:ph}{pH} 7 से ऊपर है।

\textbf{8.} $10^{7.4 - 6.1} = 10^{1.3} = 20$।

\textbf{9.} $24 / 20 = \qty{1.2}{mmol/L}$।

\textbf{10.} $10^{1.28} = 19$ और $10^{1.32} = 21$।

\textbf{11.} $20/21 = \qty{95}{\percent}$।

\textbf{12.} $10^{7.6 - 6.1} = 10^{1.5} = 32$।

\textbf{13.} हाइड्रोजनकार्बोनेट \omterm{def:g9:ions:ion}{आयन}:
\ce{HCO3- + H3O+ -> CO2 + 2H2O}।

\textbf{14.} $10^{7.1 - 6.1} = 10$।

\textbf{15.} $24/20 = 1.2$ से $24/10 = \qty{2.4}{mmol/L}$ तक: दोगुनी।
फेफड़े अतिरिक्त कार्बन डाइऑक्साइड को बाहर निकालने के लिए
तेज़ और गहरी साँस लेते हैं।

\textbf{16.} \omterm{def:g12:ka-and-pka:strong-acid}{दुर्बल अम्ल} और उसका क्षारक, दोनों में से हर एक, जो मिलाया जाए उसे सोख सकता है, और
\omterm{def:g9:acids-bases-ph:ph}{pH} केवल उनके अनुपात पर निर्भर करता है; \omterm{def:g12:ka-and-pka:strong-acid}{प्रबल अम्ल} \omterm{def:g9:acids-bases-ph:ph}{pH} को केवल
एक ही दिशा में धकेलता।

\textbf{17.} बिल्कुल नहीं: अनुपात, और इसलिए \omterm{def:g9:acids-bases-ph:ph}{pH}, वही रहते हैं।

\textbf{18.} लगभग 20।
\end{solution}
